\\maketitle
\\begin{abstract}
    本文针对\\textbf{[问题主题]}问题，建立了\\textbf{[模型类型]}模型。首先对原始数据进行\\textbf{[预处理方法]}预处理；其次构建了基于\\textbf{[核心算法]}的\\textbf{[模型名称]}模型；最后通过\\textbf{[验证方法]}验证模型的有效性。
    
    \\textbf{针对问题一}，我们建立了\\textbf{[模型1名称]}。设观测序列为$\\{y_t\\}_{t=1}^T$，将其分解为趋势项$、季节项$：
    y_t = T_t + S_t + R_t, \\quad \\sum_{t=1}^\\tau S_t = 0
    其中$\\tau。采用STL（Seasonal-Trend decomposition using Loess）方法进行稳健分解，目标函数为：
    \\min_{T, S, R} \\sum_{t=1}^T \\rho(w_t \\cdot (y_t - T_t - S_{t\\mod\\tau} - R_t))
    其中$\\rho(\\cdot)，。实验结果表明，该模型在\\textbf{[指标]}上达到\\textbf{[数值]}，优于传统方法。
    
    \\textbf{针对问题二}，我们提出了\\textbf{[模型2名称]}。假设决策变量为 \\in \\mathcal{X}$，目标函数为(x)$，约束条件为(x) \\leq 0, \\; h_k(x) = 0$。通过\\textbf{[具体方法}]优化求解：
    x^* = \\arg\\min_{x \\in \\mathcal{X}} f(x) \\quad \\text{s.t.} \\quad g_j(x) \\leq 0, \\; h_k(x) = 0
    通过对比实验，模型的\\textbf{[性能指标]}提升了\\textbf{[百分比]}\\%。
    
    \\textbf{针对问题三}，我们建立了\\textbf{[模型3名称]}。设$\\{f_i(x)\\}_{i=1}^p$，多目标优化问题可表述为：
    \\min_{x} \\; \\mathbf{f}(x) = (f_1(x), f_2(x), \\ldots, f_p(x))^T
    采用NSGA-II算法求解Pareto前沿，定义支配关系：$\\forall i, f_i(x) \\leq f_i(y)$\\exists i, f_i(x) < f_i(y)$。敏感性分析表明，模型对参数变化具有良好的鲁棒性。
    
    本文提出的\\textbf{[方法论名称]}框架，在\\textbf{[应用领域]}方面具有较好的应用价值。未来工作可考虑引入\\textbf{[改进方向]}。
    
    \\keywords{关键词1\quad 关键词2\quad 关键词3\quad 关键词4}
\\end{abstract}
